\documentclass[twocolumn]{aastex631}
\usepackage{amsmath}
\usepackage{hyperref}

\shorttitle{Adversarial Validation for Symbolic Regression}
\shortauthors{Hernandez}

\begin{document}

\title{Adversarial Validation: Stress-Testing Data-Driven Discoveries with
Automated Critique}

\author{Ivan Hernandez}
\affiliation{Independent Researcher}

\begin{abstract}
We describe a protocol for validating data-driven discoveries in astrophysics
using automated adversarial critique.  After a candidate model is discovered
via symbolic regression, an independent agent reads all code, data, and
claims, then systematically challenges every step.  The defender responds,
concedes honestly where needed, and fixes errors before the next round.
Applied across 4 projects involving Hubble constant inference, the radial
acceleration relation, black hole scaling laws, and galaxy formation
simulations, the protocol has caught 7 distinct classes of error in
pre-publication code and interpretation --- including duplicate data points,
integration bugs, column naming mismatches, degenerate statistical tests,
circular mass calibrations, and aspirational survey forecasts passed off as
realistic.  None of these errors was flagged by standard human review.
We argue that automated adversarial validation should
become a standard step in data-driven astrophysics, complementary to (not
replacing) human review.
\end{abstract}

\section{Introduction}

Data-driven methods --- symbolic regression, machine learning, automated
model discovery --- are transforming astrophysics.  They find patterns that
humans miss, scale to large datasets, and impose minimal theoretical priors.
But they also introduce new failure modes.  A symbolic regression search
might find a functional form that fits perfectly but extrapolates pathologically.
A neural network can memorize noise.  A bootstrap can conceal non-quadratic
likelihood surfaces.

Standard validation tools --- holdout tests, bootstrap, cross-validation ---
are designed to catch statistical overfitting.  They are not designed to catch
conceptual errors: mislabeled columns, circular reasoning, units confusion,
duplicate data, or aspirational claims dressed as results.

We have developed a complementary validation protocol: adversarial debate.
After the analysis is complete, an independent agent reads everything ---
data files, analysis scripts, and the manuscript --- and systematically
challenges every claim.  The defender must respond point by point,
conceding honestly.  Errors are fixed, limitations are documented, and
the process repeats until zero fatal findings remain.

This paper describes the protocol and presents 7 case studies of errors
it has caught across 4 research projects, plus a detailed worked example
from a Hubble constant analysis.

\section{Protocol}

\subsection{Initial Analysis}

The researcher completes the standard analysis pipeline:
\begin{enumerate}
  \item Parse public data into $(X,y)$ with proper errors, verified against
        published tables.
  \item Run symbolic regression (PySR~\cite{pysr}) with multiple random seeds
        and model selection by accuracy.
  \item Validate the discovered form: bootstrap resampling, holdout tests,
        cross-sample verification, literature comparison.
  \item Extend to independent data and theory predictions.
  \item Draft a manuscript reporting results and conclusions.
\end{enumerate}

\subsection{Adversarial Review}

An independent agent \textemdash{} in our implementation, a large language
model with file-reading and code-execution capabilities \textemdash{} reads
all project files and conducts a systematic audit:

\begin{itemize}
  \item \textbf{Data integrity:} Check every measurement against its
        published source.  Look for duplicates, unit mismatches, and
        transcription errors.
  \item \textbf{Methodology:} Test sensitivity to arbitrary thresholds,
        extrapolation beyond the data range, and choices of fitting
        procedure.
  \item \textbf{Interpretation:} Challenge causal claims, circular
        reasoning, and ``both methods are valid'' statements that paper
        over real discrepancies.
  \item \textbf{Presentation:} Flag overclaims, aspirational language
        masquerading as results, and missing uncertainties.
\end{itemize}

\subsection{Defender Response}

The defender responds to each challenge, categorizing it as:
\begin{itemize}
  \item \textbf{Conceded:} The error is real and must be fixed.
  \item \textbf{Partially sustained:} A genuine limitation exists but does
        not change the central result; it should be acknowledged.
  \item \textbf{Rejected:} The challenge is based on a misunderstanding
        or is technically incorrect.
\end{itemize}

Fixes are applied for all conceded and partially sustained challenges.
The resolution and any code changes are documented in a running debate log.
A subsequent round re-evaluates the fixed analysis.
The process terminates when zero fatal findings remain.

\subsection{Adversarial Prompt Template}

The following prompt is used to initialize each adversarial review round.
It is issued to a large language model with file-reading and code-execution
tools, and the model has access to the full project directory.

\begin{verbatim}
You are an adversarial reviewer auditing a research project.
Read ALL files: data, code, and the manuscript.
Then systematically challenge every claim you find. Categories:
1) Data integrity: duplicates, unit mismatches, transcription errors.
2) Methodology: sensitivity to thresholds, extrapolation, fitting choices.
3) Interpretation: circular reasoning, overclaims, hidden assumptions.
4) Presentation: missing uncertainties, aspirational language,
   undocumented priors.

For each challenge, cite specific line numbers or file paths.
Be aggressive — your job is to find problems.
The defender will respond; errors will be fixed before the next round.
\end{verbatim}

\noindent
The same prompt is re-used in subsequent rounds with the updated codebase.

\section{Case Studies}

All case studies come from projects using PySR symbolic regression applied
to astrophysical datasets.  The adversary was the opencode CLI tool using
a large language model backend.  In each case, the error had survived
standard human review before the adversarial round.

\subsection{Duplicate Data Points (BH M--$\sigma$)}

\textbf{Context:} Compilation of 126 supermassive black hole masses with
bulge velocity dispersions, drawn from multiple literature sources.

\textbf{Error:} Two independent literature searches contributed overlapping
catalogs.  The galaxy NGC 4258 appeared twice --- once from an H$_2$O maser
measurement and once from a stellar-kinematic measurement --- with different
mass values (3.9 and $3.6\times10^7\,M_\odot$) treated as separate objects.
Three other galaxies had similar duplicate entries.

\textbf{Resolution:} Deduplication logic added.  The duplicate entries
accounted for 4 of the 126 points (3\% contamination), which shifted the
best-fit slope by $0.07\pm0.04$, within uncertainty but non-negligible.

\subsection{Integration Accuracy (H$_0$ Cosmology)}

\textbf{Context:} Symbolic regression discovered an $H(z)$ polynomial from
cosmic chronometer + BAO data.  Supernova distance moduli required
numerical integration $\int_0^z dz'/H(z')$.

\textbf{Error:} The initial implementation used a Riemann sum with
$n=2000$ points at fixed redshift spacing.  At low $z$ ($<0.05$), the
$1/H(z')$ integrand varies rapidly, and 2000 points undersampled the
first $\sim$10 bins.  Error reached 0.24 mag at $z=0.01$, adding a
$7$~km/s/Mpc bias to $H_0$.

\textbf{Resolution:} Switched to Simpson's rule with adaptive spacing.
Verified convergence at $n=5000$ to $\Delta\mu < 10^{-14}$~mag.
The 7~km/s/Mpc bias was measured by re-running the full $H_0$ fit with
both integration methods and comparing the resulting $H_0$ values.

\subsection{Degenerate Statistical Null (RAR Hooks)}

\textbf{Context:} Search for non-monotonic ``hook'' features in the radial
acceleration relation (RAR) of 171 SPARC galaxies.

\textbf{Error:} The null hypothesis for hook significance was a parametric
fit (CPX5 form).  For 28 of 171 galaxies, the CPX5 fit was degenerate
($\chi^2_{\rm red} > 5$), falsely inflating apparent hook significance.
The parametric null was underfit to the data, making any residual look
significant.

\textbf{Resolution:} Replaced the parametric null with a Gaussian Process
(RBF kernel), which flexibly fits any smooth curve.  Under the GP null,
3 of 164 galaxies showed significant hooks --- consistent with the
expected 5\% false-positive rate ($p=0.99$ binomial).  The hook
population was eliminated.

\subsection{Column Name Mismatch (RAR)}

\textbf{Context:} Cross-matching simulated and observed galaxy catalogs
for RAR comparison.

\textbf{Error:} The simulation column labeled \texttt{halfmassrad} was
loaded and analyzed as the half-mass radius.  In the data pipeline,
\texttt{vmaxrad} (radius of maximum velocity, typically 2--3$\times$ larger)
had been used instead for the SPARC sample.  The mismatch propagated
through 3 analysis scripts before the adversary flagged it.

\textbf{Resolution:} Standardized on a single definition.  The original
SPARC-based result was unaffected (it used \texttt{vmaxrad} consistently),
but the simulation comparison had to be re-run with the correct definition.

\subsection{Wrong Physical Variable (Fundamental Plane)}

\textbf{Context:} Testing whether a SR-discovered galaxy scaling relation
could replace the fundamental plane for distance estimation.

\textbf{Error:} The candidate form used $M_*/R_e^2$ (stellar mass surface
density) as a proxy for surface brightness.  These quantities are
correlated but not interchangeable --- $M_*/R_e^2$ includes the
mass-to-light ratio, which varies systematically with galaxy type.
The candidate was predicting $M_*/L$ variation, not a fundamental plane.

\textbf{Resolution:} The candidate was reframed as a mass-based scaling
relation rather than a fundamental plane replacement.  The physical
interpretation changed accordingly.

\subsection{Circular Mass Calibration (BH M--$\sigma$)}

\textbf{Context:} Compilation of 126 BH masses from two methods ---
96 from stellar/gas kinematics (direct) and 30 from reverberation
mapping (RM).  RM masses are calibrated to the M--$\sigma$ relation
itself.

\textbf{Error:} The RM points were included to increase sample size,
but their masses are not independent measurements --- they assume the
M--$\sigma$ relation that the analysis was attempting to measure.
Including them in the fit biases the slope toward the input calibration.

\textbf{Resolution:} RM points flagged as not independent.  Primary
result reported from kinematic-only sample (96 points); RM-inclusive
fit reported separately as a consistency check.

\subsection{Aspirational Forecast Factors (RAR)}

\textbf{Context:} Forecast of how quickly future surveys (Euclid, Rubin,
Roman) could distinguish competing dark matter and modified gravity
models.

\textbf{Error:} The forecast used the maximum theoretical galaxy counts
from survey science requirement documents (SRDs), not realistic yields
accounting for target allocation, weather loss, and processing cuts.
For Euclid, the SRD cites $1.5\times10^9$ galaxies; realistic yields
after selection cuts are $2$--$4\times10^8$.  The factor-4 overestimate
propagated to optimistic timelines.

\textbf{Resolution:} Replaced SRD maximum counts with realistic yields
from survey operating plans.  For example, Euclid's predicted year for
distinguishing dark matter models from modified gravity shifted from
year~5 to year~7 of survey operations --- a 2--3 year delay.

\section{Worked Example: Hubble Constant}

We provide a detailed walkthrough of the adversarial protocol applied to a
joint analysis of cosmic chronometers, BAO, DESI DR2, and Pantheon+ SNe,
discovering $H(z) = H_0 + A\,z\,(z-B)(z^2+C)$ with $f(0)=0$.

The adversary raised 14 challenges across two rounds:

\subsection*{Round 1 (7 challenges)}
\begin{itemize}
  \item \textbf{r$_d$ marginalization:} BAO $r_s$ conversion uses a
        Planck prior $[146,148]$~Mpc, which does not test model-independent
        $r_s\in[130,160]$~Mpc.  \emph{Sustained:} Noted as a limitation;
        the no-DESI test (CC+SDSS only, $H_0=67.0$) provides the
        model-independent check.
  \item \textbf{H(z)-only vs joint:} $H(z)$-only fit gives $H_0=65.4$,
        $2.3\sigma$ below Planck.  \emph{Sustained:} Reflects weak
        extrapolation of the polynomial from $z>0.07$ to $z=0$; the
        SNe anchor the full joint result.
  \item \textbf{Sample-dependent $H_0$:} Three SN samples give $H_0$
        spanning 66.4--68.3 ($\sim1.9$~km/s spread).  \emph{Sustained:}
        Acknowledged as a systematic floor.
  \item \textbf{Reduced $\chi^2$:} $\chi^2_{\rm SN}/N=0.88$ suggests
        overestimated errors.  \emph{Rejected:} Conservative systematics
        are standard; the relative $\Delta\chi^2$ (fix-M test) is
        unaffected.
  \item \textbf{BAO $r_s$ circularity:} Converting $D_H/r_s$ to $H(z)$
        assumes $r_s$.  \emph{Partially sustained:} $r_s$ marginalized
        with Planck prior; no-DESI test is $r_s$-independent.
  \item \textbf{CC survey heterogeneity:} 33 points from 6 surveys.
        \emph{Partially sustained:} No cross-correlation expected,
        but noted.
  \item \textbf{Duplicate $z=0.51$:} SDSS and DESI both contribute at
        $z=0.51$.  \emph{Partially sustained:} Independent surveys;
        consistency verified ($\Delta=2.2\pm2.8$, $<1\sigma$).
\end{itemize}

\subsection*{Round 2 (7 counter-attacks)}
The second round escalated to deeper methodological challenges:
\begin{itemize}
  \item \textbf{Planck-centric r$_d$ grid} --- sustained, limitation noted.
  \item \textbf{Fix-M test symmetry} --- the test identifies inconsistency
        but not culprit; symmetry is broken by three-sample corroboration.
  \item \textbf{Bootstrap vs profile discrepancy} --- 4$\times$
        uncertainty difference explained by conditional vs.\ marginal
        inference.
  \item \textbf{Conservative covariance} --- does not affect $\Delta\chi^2$.
  \item \textbf{Coarse M(z) grid} --- adequate for $<2\sigma$ null,
        not precision.
  \item \textbf{DR1 vs DR2 mismatch in $\Lambda$CDM comparison} ---
        \emph{conceded:} $\Lambda$CDM fit used DR1 while SR used DR2;
        re-run with DR2 gave consistent results.
  \item \textbf{Duplicate paragraph in manuscript} --- \emph{conceded:}
        copy-paste error removed.
\end{itemize}

\subsection*{Outcome}
Zero fatal findings across 14 challenges.  3 conceded (technical errors),
10 partially sustained (acknowledged limitations), 1 rejected.
The core result ($H_0=68.0\pm0.9$, $>6\sigma$ SH0ES exclusion)
survived unchanged.

\section{Discussion}

\subsection{What the protocol catches}

All seven case studies come from symbolic regression (PySR) projects.
The protocol has not yet been tested on neural networks, gradient-boosted
trees, or other machine learning pipelines.  We expect the same categories
of error to appear (data integrity, circular reasoning, aspirational
forecasts) but caution that the demonstrated scope is SR-based discovery.

Across 4 projects, the 7 case studies above reveal a pattern: the errors
caught by adversarial review are \emph{conceptual}, not statistical.  None
would have been flagged by a bootstrap or cross-validation.  They arise
from:
\begin{itemize}
  \item \textbf{Propagation of initial choices} --- one wrong column name,
        one wrong integration method, one duplicated catalog entry
        silently corrupts downstream results.
  \item \textbf{Circular reasoning} --- including calibration-dependent
        data in a fit that measures the calibration.
  \item \textbf{Aspirational vs realistic} --- forecasts based on
        instrument design limits rather than operating reality.
  \item \textbf{Presentation errors} --- duplicate text, overclaimed
        significance, undocumented priors.
\end{itemize}

\subsection{What it does not catch}

The protocol is complementary to human review, not a replacement:
\begin{itemize}
  \item It cannot assess physical plausibility beyond what it has been
        trained on.
  \item It misses errors that are invisible in the code and data ---
        e.g., a systematically wrong telescope calibration.
  \item It is only as thorough as the adversary's instructions.
        A lazy adversary (``looks fine to me'') is worse than none.
\end{itemize}

\subsection{Circularity and Independence}

Because the same LLM-based tool is used for code development, manuscript
writing, and adversarial review, a concern arises: is the AI checking its
own work?  In our implementation, the adversary is invoked as a separate
session with a fresh context and explicit adversarial instructions
(Section~2.4).  It has access to all files but no memory of how the
analysis was constructed.  This independence is imperfect --- a different
LLM instance with the same underlying training data --- but the protocol
detected genuine errors even in AI-written code and text (e.g., the
duplicate paragraph in Section~4).  Users should be aware of the
residual circularity and may prefer a different model for the adversarial
role.

\subsection{Cost and Practicality}

Each adversarial round costs approximately \$2--5 in API tokens
(GPT-4-class model) and requires 30--60 minutes of wall-clock time.
The defender response takes 1--2 hours per round.  A typical project
completes in 2--3 rounds.  Of the challenges raised, approximately
40\% are rejected (adversary misunderstanding), 50\% are partially
sustained (genuine limitation, central result unaffected), and 10\%
are conceded (actual error requiring a fix).  The false-positive rate
is high, but the cost per false alarm is low --- a few minutes of
human review.

\subsection{Recommendations}

\begin{enumerate}
  \item \textbf{Make it standard:} Adversarial validation should be a
        standard step in data-driven astrophysics, as routine as
        bootstrap or cross-validation.
  \item \textbf{Document the debate:} A debate log with every challenge
        and response should accompany the paper as supplementary material.
  \item \textbf{Release the protocol:} We provide a template in our
        repository (\url{https://github.com/ivan-hernandez/h0-symbolic-regression}).
  \item \textbf{Be honest:} Concede limitations freely.  A paper that
        acknowledges 10 non-fatal limitations is more credible than one
        that claims 0.
\end{enumerate}

\section{Conclusion}

Automated adversarial validation has caught 7 distinct classes of error
across 4 projects, none of which was flagged by standard human review.  The
protocol is simple, reproducible, and complementary to existing validation
methods.  We recommend it as a standard step in data-driven astrophysics
research, particularly for projects using symbolic regression, machine
learning, or other automated discovery techniques where the conceptual
failure modes differ from traditional statistical overfitting.

\begin{acknowledgments}
This research made use of PySR~\cite{pysr}, NumPy~\cite{numpy},
Astropy~\cite{astropy}, and public data from the Pantheon+,
DES-SN5YR, DESI, and SPARC collaborations.  AI assistance (opencode)
was used for code development, adversarial review, and manuscript
preparation.
\end{acknowledgments}

\begin{thebibliography}{}
\bibitem[Cranmer(2023)]{pysr} Cranmer~M., ApJS 209, 23 (2023).
\bibitem[Harris et al.(2020)]{numpy} Harris~C.R.\ et al., Nature 585,
        357 (2020).
\bibitem[Astropy(2018)]{astropy} Astropy Collaboration, AJ 156, 123
        (2018).
\end{thebibliography}

\end{document}
